\subsection{Conclusion}

This study compared AdaBoost, XGBoost, and LightGBM for binary healthcare risk prediction across three datasets using a common experimental protocol. The protocol combined stratified development--test separation, fixed five-fold partitions, fold-local preprocessing, balanced training weights, multiple evaluation metrics, and fold-wise SHAP analysis. This design allowed predictive performance and explanation stability to be examined under the same data partitions rather than inferred from unrelated experimental settings.

For RQ1, XGBoost achieved the strongest overall ranking across the three datasets, with mean dataset rank $1.0000$, mean ROC-AUC $0.8844$, mean PR-AUC $0.5730$, and mean Recall $0.8202$. However, AdaBoost achieved the highest cross-dataset mean F1-score ($0.5392$), and the hold-out results show that the best model can depend on the dataset and metric. Therefore, the results support XGBoost as the strongest overall model under the selected ranking-oriented protocol, not as a universally superior healthcare model.

For RQ2, predictive performance was relatively stable on the two large datasets but more variable on the smaller Heart Failure Prediction dataset. This pattern was observed across the models and is consistent with the greater influence of fold composition when fewer observations are available. For RQ3, SHAP rankings were generally stable: mean Spearman correlations ranged from $0.9324$ to $0.9921$, and mean top-10 Jaccard similarity ranged from $0.7758$ to $1.0000$. The explanation results therefore support the robustness of the observed feature rankings across the evaluated partitions, while also showing that full-ranking correlation and top-feature overlap capture different aspects of stability.

The conclusions have several limitations. The study uses three public datasets and fixed, literature-guided hyperparameters rather than dataset-specific tuning. The datasets represent different prediction tasks and populations, so cross-dataset averages should be interpreted as unweighted summaries rather than evidence of clinical generalization. In addition, SHAP importance describes model behavior under the observed data distribution; it does not establish causal effects or clinical usefulness. Finally, the study evaluates stability across one five-fold partition and does not replace external validation on a new population.

\subsection{Future Work}

Future work should first test whether the stability patterns persist under repeated stratified cross-validation with multiple random partitions. This would distinguish conclusions that are robust to the fold assignment from conclusions specific to the current five-fold split. A second direction is external validation on healthcare datasets from different institutions or collection periods, together with calibration and decision-curve analysis so that ranking performance is complemented by clinically relevant probability assessment.

Further work should also analyze performance and explanation stability across demographic or clinical subgroups. This would help determine whether a high aggregate score hides weaker performance or less consistent explanations for particular populations. Finally, model-specific tuning can be studied with nested cross-validation when the research objective changes from controlled comparison to model selection. Such tuning should be evaluated separately from the current fixed-configuration experiment so that hyperparameter-selection variability is not confused with fold stability.